% Введение
\chapter*{ВВЕДЕНИЕ}
\addcontentsline{toc}{chapter}{ВВЕДЕНИЕ}

Актуальность темы курсовой работы обусловлена необходимостью автоматизации деятельности медицинской клиники, в которой ежедневно выполняются регистрация пациентов, формирование расписания врачей, запись на прием, проведение приема, расчет стоимости услуг и учет оплат. Для такой предметной области особенно важны непротиворечивость данных, защита от конфликтов по времени, сохранение истории изменений и разграничение прав доступа.

Объектом разработки является информационная система медицинской клиники. Предметом разработки являются структура базы данных, механизмы обеспечения целостности, прикладной интерфейс доступа к данным и кэширование часто используемой информации.

Целью курсовой работы является проектирование и реализация базы данных медицинской клиники и приложения доступа к ней, обеспечивающих корректную запись на прием, проведение приема, формирование счетов, учет оплат и кэширование прайс-листа и расписания врачей.

Для достижения поставленной цели необходимо решить следующие задачи:
\begin{enumerate}
    \item провести анализ предметной области медицинской клиники и существующих подходов к ее автоматизации;
    \item формализовать задачу разработки базы данных и приложения;
    \item определить сущности предметной области, связи между ними и ограничения целостности;
    \item спроектировать ролевую модель базы данных, процедуры, триггеры и транзакционные сценарии;
    \item выбрать средства реализации базы данных, backend- и frontend-частей приложения;
    \item реализовать структуру базы данных, серверную часть, кэширование в Redis и демонстрационный пользовательский интерфейс;
    \item разработать тесты для проверки записи, отмены, проведения приема, формирования счетов и регистрации оплаты;
    \item исследовать конкурентный доступ и влияние индексов на производительность ключевых операций.
\end{enumerate}

\clearpage
